% Recommended: compile with XeLaTeX or LuaLaTeX.
% pdfLaTeX is also supported through the fallback preamble below.
\documentclass[12pt,a4paper]{report}

\usepackage[a4paper,margin=2.5cm]{geometry}
\usepackage{iftex}
\ifPDFTeX
    \usepackage[T5]{fontenc}
    \usepackage[utf8]{inputenc}
    \usepackage[vietnamese]{babel}
\else
    \usepackage{fontspec}
    \usepackage{polyglossia}
    \setdefaultlanguage{vietnamese}
    \defaultfontfeatures{Ligatures=TeX}
    \setmainfont{TeX Gyre Termes}
\fi
\usepackage{setspace}
\usepackage{hyperref}
\usepackage{graphicx}
\usepackage{booktabs}
\usepackage{longtable}
\usepackage{array}
\usepackage{float}
\usepackage{enumitem}
\usepackage{xcolor}
\usepackage{tikz}
\usetikzlibrary{positioning, arrows.meta, calc}

\hypersetup{
    colorlinks=true,
    linkcolor=blue,
    urlcolor=blue,
    pdftitle={Bao cao database he thong Pody},
    pdfauthor={Nhom sinh vien Pody}
}

\setstretch{1.25}
\setlength{\parskip}{0.5em}
\setlength{\parindent}{0pt}

\newcommand{\diagramplaceholder}[2]{
\begin{figure}[H]
\centering
\fbox{
\parbox{0.88\textwidth}{
\centering
\vspace{1.0cm}
#1\\[0.5em]
\textit{Chèn sơ đồ Mermaid đã dựng trước đó tại đây sau khi xuất sang ảnh hoặc PDF.}
\vspace{1.0cm}
}}
\caption{#2}
\end{figure}
}

\tikzset{
    servicebox/.style={
        draw,
        rounded corners=4pt,
        thick,
        fill=green!8,
        align=center,
        font=\bfseries\small,
        minimum width=3.2cm,
        minimum height=1cm
    },
    datastorebox/.style={
        draw,
        rounded corners=4pt,
        thick,
        fill=orange!10,
        align=center,
        font=\small,
        minimum width=3.2cm,
        minimum height=1cm
    },
    entity/.style={
        draw,
        rounded corners=3pt,
        thick,
        fill=blue!4,
        align=left,
        font=\ttfamily\scriptsize,
        text width=3.8cm,
        inner sep=5pt
    },
    entitywide/.style={
        entity,
        text width=4.4cm
    },
    relation/.style={
        -{Latex[length=2.2mm]},
        thick
    },
    softrelation/.style={
        -{Latex[length=2.2mm]},
        thick,
        dashed
    },
    relabel/.style={
        font=\scriptsize\itshape,
        fill=white,
        inner sep=1pt
    }
}

\newcommand{\entitytitle}[1]{\textbf{\textsf{#1}}\\\rule{\linewidth}{0.3pt}}

\newcommand{\systemdbdiagram}{
\begin{figure}[H]
\centering
\resizebox{\textwidth}{!}{
\begin{tikzpicture}[node distance=1.4cm and 2.0cm]
    \node[servicebox] (identity) {Identity-Service};
    \node[servicebox, below=of identity] (content) {Content-Service};
    \node[servicebox, below=of content] (ai) {AI-Service};
    \node[servicebox, below=of ai] (article) {Article-Service};
    \node[servicebox, below=of article] (notification) {Notification-Service};
    \node[servicebox, below=of notification] (billing) {Billing-Service};
    \node[servicebox, below=of billing] (social) {Social-Service};
    \node[servicebox, below=of social] (embedding) {Embedding-Service};

    \node[datastorebox, right=of identity] (identitydb) {PostgreSQL\\Identity DB};
    \node[datastorebox, right=of content] (contentdb) {PostgreSQL\\Content DB};
    \node[datastorebox, right=of ai] (aidb) {PostgreSQL\\AI DB};
    \node[datastorebox, right=of article] (articledb) {PostgreSQL\\Article DB};
    \node[datastorebox, right=of notification] (notificationdb) {PostgreSQL\\Notification DB};
    \node[datastorebox, right=of billing] (billingdb) {PostgreSQL\\Billing DB\\(thiết kế)};
    \node[datastorebox, right=of social] (socialdb) {MongoDB\\Social DB\\(thiết kế)};
    \node[datastorebox, right=of embedding] (embeddingdb) {PostgreSQL + pgvector\\Embedding DB\\(đề xuất)};

    \foreach \a/\b in {
        identity/identitydb,
        content/contentdb,
        ai/aidb,
        article/articledb,
        notification/notificationdb,
        billing/billingdb,
        social/socialdb,
        embedding/embeddingdb
    }{
        \draw[relation] (\a) -- (\b);
    }

    \node[draw, rounded corners=4pt, thick, fill=gray!10, align=center, font=\small, minimum width=4.3cm, minimum height=1cm, below=1.8cm of embeddingdb] (bus) {Event Bus / External ID};

    \foreach \srv in {identity,content,ai,article,notification,billing,social,embedding}{
        \draw[softrelation] (\srv.east) |- (bus.west);
    }
\end{tikzpicture}
}
\caption{Sơ đồ tổng quan cơ sở dữ liệu toàn hệ thống}
\end{figure}
}

\newcommand{\contentservicediagram}{
\begin{figure}[H]
\centering
\resizebox{\textwidth}{!}{
\begin{tikzpicture}[node distance=1.1cm and 1.4cm]
    \node[entity] (categories) {\entitytitle{CATEGORIES}id PK\\slug UK\\name\\applies\_to\\sort\_order};
    \node[entity, right=of categories] (showcategories) {\entitytitle{SHOW\_CATEGORIES}show\_id PK, FK\\category\_id PK, FK\\is\_primary};
    \node[entitywide, right=of showcategories] (shows) {\entitytitle{SHOWS}id PK\\owner\_user\_id\\title\\slug UK\\content\_type\\publish\_status\\visibility};
    \node[entity, right=of shows] (showtags) {\entitytitle{SHOW\_TAGS}show\_id PK, FK\\tag\_id PK, FK};
    \node[entity, right=of showtags] (tags) {\entitytitle{TAGS}id PK\\slug UK\\name};

    \node[entity, below=of shows] (showhosts) {\entitytitle{SHOW\_HOSTS}id PK\\show\_id FK\\display\_name\\role\\persona\_type\\sort\_order};
    \node[entitywide, below=of showhosts] (episodes) {\entitytitle{EPISODES}id PK\\show\_id FK\\title\\slug\\season\_number\\episode\_number\\audio\_url\\publish\_status};
    \node[entity, right=of episodes] (episodeassets) {\entitytitle{EPISODE\_ASSETS}id PK\\episode\_id FK\\asset\_type\\url\\storage\_key};
    \node[entity, left=of episodes] (episodesegments) {\entitytitle{EPISODE\_SEGMENTS}id PK\\episode\_id FK\\show\_host\_id FK\\segment\_index\\text\_content};
    \node[entity, below=of episodes] (blocks) {\entitytitle{EPISODE\_COMPANION\_BLOCKS}id PK\\episode\_id FK\\block\_type\\sort\_order\\payload};

    \draw[relation] (categories) -- node[relabel, above] {used in} (showcategories);
    \draw[relation] (showcategories) -- node[relabel, above] {classifies} (shows);
    \draw[relation] (shows) -- node[relabel, above] {tagged} (showtags);
    \draw[relation] (showtags) -- node[relabel, above] {uses} (tags);
    \draw[relation] (shows) -- node[relabel, right] {has} (showhosts);
    \draw[relation] (showhosts) -- node[relabel, right] {contains} (episodes);
    \draw[relation] (episodes) -- node[relabel, above] {has} (episodeassets);
    \draw[relation] (episodes) -- node[relabel, above] {has} (episodesegments);
    \draw[relation] (episodes) -- node[relabel, right] {has} (blocks);
    \draw[softrelation] (showhosts) -- node[relabel, above] {speaker} (episodesegments);
\end{tikzpicture}
}
\caption{Sơ đồ quan hệ dữ liệu của Content-Service}
\end{figure}
}

\newcommand{\aiservicediagram}{
\begin{figure}[H]
\centering
\resizebox{\textwidth}{!}{
\begin{tikzpicture}[node distance=1.1cm and 1.4cm]
    \node[entity] (threads) {\entitytitle{CHAT\_THREADS}id PK\\owner\_user\_id\\title\\status};
    \node[entity, below=of threads] (messages) {\entitytitle{CHAT\_MESSAGES}id PK\\thread\_id FK\\role\\text\_content\\plan\_id FK};
    \node[entity, below=of messages] (attachments) {\entitytitle{CHAT\_ATTACHMENTS}id PK\\message\_id FK\\attachment\_type\\file\_name\\url};

    \node[entitywide, right=of threads] (plans) {\entitytitle{PRODUCTION\_PLANS}id PK\\owner\_user\_id\\thread\_id FK\\target\_show\_id\\series\_title\\status\\target\_language\_code};
    \node[entity, right=of plans] (voices) {\entitytitle{VOICE\_PROFILES}id PK\\name\\provider\\provider\_voice\_id\\language\_code};
    \node[entity, below=of plans] (hosts) {\entitytitle{PLAN\_HOSTS}id PK\\plan\_id FK\\voice\_profile\_id FK\\display\_name\\role};
    \node[entity, below left=of hosts] (sources) {\entitytitle{PLAN\_SOURCES}id PK\\plan\_id FK\\source\_article\_id\\source\_title\\source\_url};
    \node[entity, below right=of hosts] (drafts) {\entitytitle{EPISODE\_DRAFTS}id PK\\plan\_id FK\\generated\_episode\_id\\episode\_number\\title\\status};
    \node[entity, below=of drafts] (jobs) {\entitytitle{GENERATION\_JOBS}id PK\\plan\_id FK\\episode\_draft\_id FK\\job\_type\\status\\provider};

    \draw[relation] (threads) -- node[relabel, right] {has} (messages);
    \draw[relation] (messages) -- node[relabel, right] {has} (attachments);
    \draw[relation] (threads) -- node[relabel, above] {backs} (plans);
    \draw[softrelation] (messages) -- node[relabel, above] {references} (plans);
    \draw[relation] (plans) -- node[relabel, right] {has} (hosts);
    \draw[relation] (voices) -- node[relabel, above] {voices} (hosts);
    \draw[relation] (plans) -- node[relabel, above] {uses} (sources);
    \draw[relation] (plans) -- node[relabel, above] {drafts} (drafts);
    \draw[relation] (plans) -- node[relabel, right] {runs} (jobs);
    \draw[softrelation] (drafts) -- node[relabel, right] {feeds} (jobs);
\end{tikzpicture}
}
\caption{Sơ đồ quan hệ dữ liệu của AI-Service}
\end{figure}
}

\newcommand{\articleservicediagram}{
\begin{figure}[H]
\centering
\resizebox{0.96\textwidth}{!}{
\begin{tikzpicture}[node distance=1.2cm and 1.5cm]
    \node[entity] (sources) {\entitytitle{NEWS\_SOURCES}id PK\\domain UK\\name\\rss\_url UK\\status};
    \node[entitywide, right=of sources] (articles) {\entitytitle{ARTICLES}id PK\\source\_id\\title\\summary\\content\\original\_url UK\\published\_at};
    \node[entity, above right=of articles] (categories) {\entitytitle{ARTICLE\_CATEGORIES}id PK\\article\_id\\category\_name};
    \node[entity, below right=of articles] (stats) {\entitytitle{ARTICLE\_STATS}article\_id PK\\view\_count\\updated\_at};
    \node[entity, below=of articles] (interactions) {\entitytitle{ARTICLE\_INTERACTIONS}id PK\\article\_id\\user\_id\\interaction\_type};
    \node[entity, below left=of articles] (metrics) {\entitytitle{ARTICLE\_METRICS}id PK\\article\_id\\user\_id\\reading\_time\_seconds};
    \node[entity, above left=of articles] (comments) {\entitytitle{ARTICLE\_COMMENTS}id PK\\article\_id\\user\_id\\user\_name\\content};

    \draw[relation] (sources) -- node[relabel, above] {publishes} (articles);
    \draw[relation] (articles) -- node[relabel, above] {categorized} (categories);
    \draw[relation] (articles) -- node[relabel, right] {has stats} (stats);
    \draw[relation] (articles) -- node[relabel, right] {receives} (interactions);
    \draw[relation] (articles) -- node[relabel, above] {measured by} (metrics);
    \draw[relation] (articles) -- node[relabel, above] {commented by} (comments);
\end{tikzpicture}
}
\caption{Sơ đồ quan hệ dữ liệu của Article-Service}
\end{figure}
}

\newcommand{\embeddingservicediagram}{
\begin{figure}[H]
\centering
\resizebox{\textwidth}{!}{
\begin{tikzpicture}[node distance=1.2cm and 1.5cm]
    \node[entity] (models) {\entitytitle{EMBEDDING\_MODELS}id PK\\code UK\\provider\\model\_name\\dimensions};
    \node[entitywide, right=of models] (docs) {\entitytitle{ARTICLE\_EMBEDDING\_DOCUMENTS}id PK\\article\_id UK\\source\_id\\title\\summary\\content\_hash};
    \node[entity, below=of docs] (chunks) {\entitytitle{ARTICLE\_EMBEDDING\_CHUNKS}id PK\\document\_id FK\\chunk\_index\\chunk\_type\\content};
    \node[entity, right=of chunks] (chunkemb) {\entitytitle{ARTICLE\_CHUNK\_EMBEDDINGS}id PK\\chunk\_id FK\\model\_id FK\\embedding\\embedding\_version};
    \node[entity, above right=of chunkemb] (catdocs) {\entitytitle{CATEGORY\_EMBEDDING\_DOCUMENTS}id PK\\category\_name UK\\display\_name\\description\\keywords};
    \node[entity, below=of catdocs] (catemb) {\entitytitle{CATEGORY\_EMBEDDINGS}id PK\\category\_doc\_id FK\\model\_id FK\\embedding};
    \node[entity, below left=of chunkemb] (jobs) {\entitytitle{EMBEDDING\_JOBS}id PK\\target\_type\\target\_id\\model\_id FK\\job\_type\\status};

    \draw[relation] (models) -- node[relabel, above] {uses} (chunkemb);
    \draw[relation] (models) -- node[relabel, above] {uses} (catemb);
    \draw[relation] (models) -- node[relabel, left] {requested by} (jobs);
    \draw[relation] (docs) -- node[relabel, right] {split into} (chunks);
    \draw[relation] (chunks) -- node[relabel, above] {embedded as} (chunkemb);
    \draw[relation] (catdocs) -- node[relabel, right] {embedded as} (catemb);
\end{tikzpicture}
}
\caption{Sơ đồ quan hệ dữ liệu của Embedding-Service}
\end{figure}
}

\newcommand{\identityservicediagram}{
\begin{figure}[H]
\centering
\resizebox{\textwidth}{!}{
\begin{tikzpicture}[node distance=1.2cm and 1.5cm]
    \node[entitywide] (users) {\entitytitle{USERS}id PK\\email UK\\display\_name\\username UK\\account\_type\\status\\email\_verified\_at};
    \node[entity, above left=of users] (identities) {\entitytitle{USER\_IDENTITIES}id PK\\user\_id FK\\provider\\provider\_user\_id\\provider\_email};
    \node[entity, above right=of users] (devices) {\entitytitle{USER\_DEVICES}id PK\\user\_id FK\\platform\\push\_provider\\device\_token UK};
    \node[entity, below right=of users] (sessions) {\entitytitle{AUTH\_SESSIONS}id PK\\user\_id FK\\user\_device\_id FK\\refresh\_token\_hash UK\\expires\_at};
    \node[entity, below left=of users] (emailtokens) {\entitytitle{EMAIL\_VERIFICATION\_TOKENS}id PK\\user\_id FK\\token\_hash UK\\expires\_at\\used\_at};
    \node[entity, below=of users] (resettokens) {\entitytitle{PASSWORD\_RESET\_TOKENS}id PK\\user\_id FK\\token\_hash UK\\expires\_at\\used\_at};

    \draw[relation] (users) -- node[relabel, above] {has} (identities);
    \draw[relation] (users) -- node[relabel, above] {has} (devices);
    \draw[relation] (users) -- node[relabel, right] {has} (sessions);
    \draw[relation] (users) -- node[relabel, left] {verifies} (emailtokens);
    \draw[relation] (users) -- node[relabel, right] {resets} (resettokens);
    \draw[softrelation] (devices) -- node[relabel, right] {used by} (sessions);
\end{tikzpicture}
}
\caption{Sơ đồ quan hệ dữ liệu của Identity-Service}
\end{figure}
}

\newcommand{\notificationservicediagram}{
\begin{figure}[H]
\centering
\resizebox{0.95\textwidth}{!}{
\begin{tikzpicture}[node distance=1.3cm and 1.7cm]
    \node[entity] (settings) {\entitytitle{USER\_NOTIFICATION\_SETTINGS}user\_id PK\\push\_enabled\\email\_enabled\\new\_episode\_enabled\\comment\_enabled};
    \node[entitywide, right=of settings] (notifications) {\entitytitle{NOTIFICATIONS}id PK\\user\_id\\actor\_user\_id\\type\\target\_type\\target\_id\\title\\is\_read};
    \node[entity, right=of notifications] (logs) {\entitytitle{DELIVERY\_LOGS}id PK\\notification\_id FK\\user\_id\\channel\\provider\\delivery\_status};
    \node[entity, below=of notifications] (templates) {\entitytitle{TEMPLATES}id PK\\template\_code\\channel\\title\_template\\body\_template\\is\_active};

    \draw[relation] (notifications) -- node[relabel, above] {has} (logs);
    \draw[softrelation] (settings) -- node[relabel, above] {configures} (notifications);
    \draw[softrelation] (templates) -- node[relabel, right] {formats} (notifications);
\end{tikzpicture}
}
\caption{Sơ đồ quan hệ dữ liệu của Notification-Service}
\end{figure}
}

\newcommand{\billingservicediagram}{
\begin{figure}[H]
\centering
\resizebox{\textwidth}{!}{
\begin{tikzpicture}[node distance=1.2cm and 1.5cm]
    \node[entity] (plans) {\entitytitle{SUBSCRIPTION\_PLANS}id PK\\code UK\\name\\billing\_period\\numeric12\_2 price\_amount};
    \node[entitywide, right=of plans] (subs) {\entitytitle{USER\_SUBSCRIPTIONS}id PK\\user\_id\\plan\_id FK\\provider\\provider\_subscription\_id\\status};
    \node[entity, below=of subs] (payments) {\entitytitle{PAYMENT\_TRANSACTIONS}id PK\\user\_id\\subscription\_id FK\\transaction\_type\\status\\numeric12\_2 amount};
    \node[entity, below=of plans] (ledger) {\entitytitle{CREDIT\_LEDGER}id PK\\user\_id\\delta\_credits\\balance\_after\\entry\_type\\reference\_type};
    \node[entity, below=of ledger] (reservations) {\entitytitle{CREDIT\_RESERVATIONS}id PK\\user\_id\\purpose\\reference\_type\\reserved\_credits\\status};
    \node[entity, right=of ledger] (entitlements) {\entitytitle{CONTENT\_ENTITLEMENTS}id PK\\user\_id\\show\_id\\episode\_id\\access\_type\\granted\_by\_ledger\_id FK};

    \draw[relation] (plans) -- node[relabel, above] {selected by} (subs);
    \draw[relation] (subs) -- node[relabel, right] {billed in} (payments);
    \draw[relation] (ledger) -- node[relabel, above] {grants} (entitlements);
    \draw[softrelation] (ledger) -- node[relabel, left] {reserves} (reservations);
\end{tikzpicture}
}
\caption{Sơ đồ quan hệ dữ liệu của Billing-Service}
\end{figure}
}

\newcommand{\socialservicediagram}{
\begin{figure}[H]
\centering
\resizebox{\textwidth}{!}{
\begin{tikzpicture}[node distance=1.2cm and 1.4cm]
    \node[entity] (userfollows) {\entitytitle{USER\_FOLLOWS}\_id PK\\follower\_user\_id\\followed\_user\_id\\follower\_snapshot\\followed\_snapshot};
    \node[entity, right=of userfollows] (showfollows) {\entitytitle{SHOW\_FOLLOWS}\_id PK\\user\_id\\show\_id\\notifications\_enabled\\show\_snapshot};
    \node[entity, right=of showfollows] (saved) {\entitytitle{SAVED\_EPISODES}\_id PK\\user\_id\\episode\_id\\show\_id\\episode\_snapshot};

    \node[entity, below=of userfollows] (playlists) {\entitytitle{PLAYLISTS}\_id PK\\owner\_user\_id\\name\\is\_public\\item\_count};
    \node[entity, right=of playlists] (playlistitems) {\entitytitle{PLAYLIST\_ITEMS}\_id PK\\playlist\_id\\episode\_id\\show\_id\\sort\_order};
    \node[entity, right=of playlistitems] (progress) {\entitytitle{LISTENING\_PROGRESS}\_id PK\\user\_id\\episode\_id\\position\_seconds\\progress\_ratio};

    \node[entity, below=of playlists] (reactions) {\entitytitle{EPISODE\_REACTIONS}\_id PK\\episode\_id\\user\_id\\reaction\_type\\user\_snapshot};
    \node[entity, right=of reactions] (comments) {\entitytitle{COMMENTS}\_id PK\\episode\_id\\user\_id\\parent\_comment\_id\\body\\user\_snapshot};
    \node[entity, right=of comments] (counters) {\entitytitle{ENGAGEMENT\_COUNTERS}\_id PK\\episode\_id\\like\_count\\comment\_count\\save\_count};

    \draw[softrelation] (playlists) -- node[relabel, above] {contains} (playlistitems);
    \draw[softrelation] (reactions) -- node[relabel, above] {updates} (counters);
    \draw[softrelation] (comments) -- node[relabel, above] {updates} (counters);
    \draw[softrelation] (saved) -- node[relabel, left] {feeds} (counters);
    \draw[softrelation] (showfollows) -- node[relabel, left] {belongs to user} (progress);
\end{tikzpicture}
}
\caption{Sơ đồ logic dữ liệu của Social-Service}
\end{figure}
}

\begin{document}

\begin{titlepage}
\centering

{\Large \textbf{BÁO CÁO THIẾT KẾ CƠ SỞ DỮ LIỆU HỆ THỐNG PODY}\par}
\vspace{1cm}
{\large Phục vụ mô tả kiến trúc dữ liệu và phân công cá nhân trong nhóm\par}
\vspace{2cm}

\begin{tabular}{|p{5.2cm}|p{4.2cm}|p{4cm}|}
\hline
\textbf{Họ và tên} & \textbf{Mã sinh viên} & \textbf{Phụ trách} \\
\hline
Lại Xuân Hiếu & B22DCCN309 & Content-Service, AI-Service \\
\hline
Vũ Đức Thành & B22DCCN801 & Article-Service, Embedding-Service \\
\hline
Nguyễn Việt Huy & B22DCCN393 & Identity-Service, Notification-Service \\
\hline
Lê Minh Ngọc & 586 & Billing-Service, Social-Service \\
\hline
\end{tabular}

\vfill

{\large Nhóm thực hiện: 04 thành viên\par}
\vspace{0.5cm}
{\large \today\par}

\end{titlepage}

\tableofcontents
\clearpage

\chapter{Mô tả database toàn hệ thống}

\section{Tổng quan kiến trúc dữ liệu}

Hệ thống Pody được xây dựng theo kiến trúc microservices, trong đó mỗi service sở hữu một không gian dữ liệu riêng. Cách tổ chức này giúp từng service có thể phát triển độc lập, giảm mức độ phụ thuộc lẫn nhau và thuận lợi hơn trong việc triển khai, mở rộng cũng như bảo trì. Về bản chất, thiết kế database của hệ thống không đi theo mô hình một cơ sở dữ liệu tập trung, mà đi theo hướng phân tách theo domain nghiệp vụ.

Trong phạm vi hiện tại của dự án, các service quan trọng như \texttt{Identity-Service}, \texttt{Content-Service}, \texttt{AI-Service}, \texttt{Article-Service} và \texttt{Notification-Service} đang sử dụng PostgreSQL. Ngoài ra, hệ thống còn có các phần thiết kế mở rộng cho \texttt{Embedding-Service}, \texttt{Billing-Service} và \texttt{Social-Service}. Trong đó, \texttt{Billing-Service} đã có schema tương đối đầy đủ, còn \texttt{Embedding-Service} và \texttt{Social-Service} phản ánh định hướng mở rộng cho semantic search, recommendation và tương tác người dùng.

\section{Nguyên tắc thiết kế cơ sở dữ liệu}

Thiết kế cơ sở dữ liệu của hệ thống dựa trên các nguyên tắc chính sau đây:

\begin{enumerate}[label=\arabic*.]
    \item \textbf{Database per service}: mỗi service quản lý schema và dữ liệu của chính nó, không chia sẻ bảng trực tiếp với service khác.
    \item \textbf{Không có foreign key xuyên service}: các service chỉ lưu external id để tham chiếu, thay vì ràng buộc vật lý giữa các database.
    \item \textbf{Đồng bộ bất đồng bộ qua event}: dữ liệu giữa các service được cập nhật thông qua event bus thay vì truy vấn chéo trực tiếp.
    \item \textbf{Áp dụng mẫu outbox/inbox}: nhiều service có cặp bảng \texttt{outbox\_events} và \texttt{inbox\_processed\_events} để hỗ trợ phát event ổn định và xử lý idempotent.
    \item \textbf{Tách dữ liệu lớn ra khỏi database}: audio, ảnh, transcript và các tệp lớn không lưu trực tiếp trong database mà lưu ở object storage; database chỉ lưu URL hoặc metadata.
\end{enumerate}

\section{Phân loại các kho lưu trữ trong hệ thống}

\begin{longtable}{|p{3.8cm}|p{3.3cm}|p{3.1cm}|p{4.6cm}|}
\hline
\textbf{Service} & \textbf{Kho lưu trữ chính} & \textbf{Trạng thái} & \textbf{Ghi chú} \\
\hline
Identity-Service & PostgreSQL & Đang dùng & Quản lý user, provider đăng nhập, session, device \\
\hline
Notification-Service & PostgreSQL & Đang dùng & Lưu thông báo, cấu hình thông báo, lịch sử gửi \\
\hline
Content-Service & PostgreSQL & Đang dùng & Lưu show, host, episode, asset, segment và companion block \\
\hline
AI-Service & PostgreSQL & Đang dùng & Lưu thread, tin nhắn, production plan, draft, generation job \\
\hline
Article-Service & PostgreSQL & Đang dùng & Lưu nguồn tin, bài báo, category, comment, interaction, metric \\
\hline
Embedding-Service & PostgreSQL + pgvector & Đề xuất mở rộng & Phục vụ semantic search và vector retrieval \\
\hline
Billing-Service & PostgreSQL & Có schema, chưa thấy wire runtime & Quản lý subscription, transaction, credit, entitlement \\
\hline
Social-Service & MongoDB theo thiết kế & Mức thiết kế & Quản lý follow, reaction, comment, playlist, listening progress \\
\hline
\end{longtable}

\section{Mối quan hệ dữ liệu ở mức toàn hệ thống}

Ở mức hệ thống, mỗi service quản lý một domain riêng:

\begin{itemize}
    \item \textbf{Identity-Service} quản lý danh tính người dùng và thông tin xác thực.
    \item \textbf{Content-Service} quản lý nội dung cốt lõi của nền tảng như show, host và episode.
    \item \textbf{AI-Service} phục vụ quá trình lên kế hoạch và sinh nội dung bằng AI.
    \item \textbf{Article-Service} quản lý bài báo, nguồn tin và dữ liệu tương tác với bài báo.
    \item \textbf{Notification-Service} quản lý thông báo trong ứng dụng và log gửi thông báo.
    \item \textbf{Billing-Service} quản lý giao dịch, subscription, credit và entitlement.
    \item \textbf{Social-Service} quản lý hành vi xã hội như follow, reaction, comment và playlist.
\end{itemize}

Tuy các service có liên hệ nghiệp vụ với nhau, nhưng mối liên hệ này chủ yếu được thể hiện thông qua external id hoặc event. Ví dụ:

\begin{itemize}
    \item \texttt{Content-Service} có thể lưu \texttt{owner\_user\_id} để trỏ đến người dùng trong \texttt{Identity-Service}.
    \item \texttt{AI-Service} có thể lưu \texttt{target\_show\_id} hoặc \texttt{source\_article\_id} để liên kết logic đến \texttt{Content-Service} và \texttt{Article-Service}.
    \item \texttt{Billing-Service} lưu \texttt{user\_id}, \texttt{show\_id}, \texttt{episode\_id} như external id thay vì foreign key vật lý.
\end{itemize}

Nhờ vậy, các database vẫn tách biệt nhưng toàn hệ thống vẫn giữ được tính liên thông ở tầng nghiệp vụ.

\section{Đánh giá tổng quan về thiết kế database}

Thiết kế cơ sở dữ liệu của Pody có một số ưu điểm nổi bật:

\begin{itemize}
    \item tách biệt rõ ràng giữa các domain nghiệp vụ;
    \item dễ mở rộng từng service mà không tác động trực tiếp lên service khác;
    \item phù hợp với kiến trúc event-driven;
    \item thuận lợi cho việc mở rộng thêm các service đặc thù như Embedding hoặc Billing.
\end{itemize}

Bên cạnh đó, mô hình này cũng đặt ra một số yêu cầu:

\begin{itemize}
    \item phải quản lý tốt luồng đồng bộ dữ liệu giữa các service;
    \item cần cẩn thận với tính nhất quán cuối cùng giữa các read model;
    \item phải mô tả rõ external id để tránh nhầm lẫn khi tích hợp.
\end{itemize}

\systemdbdiagram

\chapter{Mô tả cá nhân}

\section{Lại Xuân Hiếu -- MSSV: B22DCCN309}

\subsection{Content-Service}

Content-Service là service trọng tâm trong hệ thống vì đây là nơi lưu trữ nội dung chính mà người dùng tiêu thụ. Database của service này được thiết kế theo hướng quan hệ chặt chẽ, trong đó \texttt{shows} và \texttt{episodes} là hai bảng lõi.

Từ góc nhìn cơ sở dữ liệu, \texttt{shows} đại diện cho các chương trình hoặc series, còn \texttt{episodes} đại diện cho từng tập cụ thể. Một show có thể có nhiều host, nhiều tag, nhiều category và nhiều episode. Mỗi episode lại có thể có transcript segment, asset đa phương tiện và companion block phục vụ hiển thị nội dung nâng cao trong ứng dụng.

Các bảng quan trọng nhất của service này gồm:

\begin{itemize}
    \item \texttt{shows}: lưu thông tin tổng quát của một show như tiêu đề, mô tả, ảnh bìa, trạng thái phát hành và thông tin chủ sở hữu.
    \item \texttt{show\_hosts}: lưu danh sách host gắn với show.
    \item \texttt{episodes}: lưu thông tin từng tập, bao gồm audio URL, thứ tự tập, thời lượng và trạng thái phát hành.
    \item \texttt{episode\_segments}: lưu transcript được tách đoạn theo speaker.
    \item \texttt{episode\_assets}: lưu metadata về các tài nguyên đi kèm như ảnh, subtitle hoặc transcript.
    \item \texttt{episode\_companion\_blocks}: lưu các khối nội dung mở rộng như timeline, flashcards, quote hoặc quiz.
\end{itemize}

Điểm mạnh của thiết kế này là khả năng biểu diễn rất rõ quan hệ cha -- con giữa show và episode, đồng thời hỗ trợ tốt cho các chức năng hiển thị nội dung phong phú trong ứng dụng. Việc tách category và tag thành bảng riêng cùng các bảng liên kết trung gian giúp hệ thống dễ mở rộng bộ lọc, tìm kiếm và phân loại nội dung.

\contentservicediagram

\subsection{AI-Service}

AI-Service là service chịu trách nhiệm hỗ trợ người dùng lập kế hoạch và sinh nội dung bằng trí tuệ nhân tạo. Nếu Content-Service lưu sản phẩm cuối cùng, thì AI-Service lưu toàn bộ quá trình chuẩn bị, tạo bản nháp và quản lý job sinh nội dung.

Database của AI-Service có thể chia thành bốn lớp:

\begin{itemize}
    \item lớp hội thoại: \texttt{chat\_threads}, \texttt{chat\_messages}, \texttt{chat\_attachments};
    \item lớp cấu hình giọng nói: \texttt{voice\_profiles};
    \item lớp kế hoạch sản xuất: \texttt{production\_plans}, \texttt{production\_plan\_hosts}, \texttt{production\_plan\_sources}, \texttt{production\_plan\_episode\_drafts};
    \item lớp thực thi: \texttt{generation\_jobs}.
\end{itemize}

Bảng \texttt{production\_plans} đóng vai trò trung tâm, liên kết đến thread hội thoại, danh sách host, nguồn tư liệu và các bản nháp episode. Từ đó, \texttt{generation\_jobs} theo dõi trạng thái chạy của từng tác vụ như sinh script, sinh audio, sinh ảnh hoặc publish episode.

Thiết kế này đặc biệt phù hợp cho workflow AI vì nó cho phép:

\begin{itemize}
    \item lưu lịch sử trao đổi với người dùng;
    \item quản lý trạng thái của từng bản kế hoạch;
    \item liên kết giữa plan và dữ liệu đầu vào như bài báo hoặc voice profile;
    \item theo dõi vòng đời job từ lúc tạo đến lúc hoàn tất hoặc thất bại.
\end{itemize}

\aiservicediagram

\section{Vũ Đức Thành -- MSSV: B22DCCN801}

\subsection{Article-Service}

Article-Service chịu trách nhiệm quản lý toàn bộ dữ liệu liên quan đến bài báo và nguồn tin. Đây là service nền tảng cho khối nội dung thời sự, phân tích và các tính năng AI dựa trên thông tin bài báo.

Database của service này xoay quanh hai đối tượng trung tâm là nguồn tin và bài báo. Bảng \texttt{news\_sources} lưu thông tin về domain, tên nguồn, RSS URL và độ tin cậy. Bảng \texttt{articles} lưu bài báo thực tế với tiêu đề, tóm tắt, nội dung, URL gốc và thời điểm phát hành.

Ngoài ra, service còn có các bảng bổ sung nhằm hỗ trợ hiển thị và tương tác:

\begin{itemize}
    \item \texttt{article\_categories}: gắn thể loại cho bài báo;
    \item \texttt{article\_stats}: thống kê lượt xem;
    \item \texttt{article\_interactions}: lưu tương tác như like hoặc love;
    \item \texttt{article\_metrics}: lưu dữ liệu về thời gian đọc;
    \item \texttt{article\_comments}: lưu bình luận của người dùng.
\end{itemize}

Điểm đáng chú ý là nhiều quan hệ trong service này hiện được duy trì ở mức logic trong tầng ứng dụng, thay vì foreign key vật lý trong SQL. Điều đó giúp service linh hoạt hơn trong giai đoạn phát triển, nhưng cũng yêu cầu tầng repository và service phải kiểm soát dữ liệu chặt chẽ hơn.

\articleservicediagram

\subsection{Embedding-Service}

Embedding-Service là phần mở rộng có vai trò đặc biệt trong hệ thống, phục vụ các bài toán semantic search, vector retrieval và hỗ trợ AI truy xuất ngữ nghĩa từ bài báo hoặc thể loại báo. Dù chưa có migration chính thức trong repo, đây là hướng thiết kế hợp lý để nâng cấp hệ thống.

Khác với Article-Service vốn lưu dữ liệu gốc, Embedding-Service tập trung vào dữ liệu đã được chuẩn hóa để sinh vector. Thiết kế được chia thành các lớp:

\begin{itemize}
    \item cấu hình model: \texttt{embedding\_models};
    \item tài liệu bài báo và chunk: \texttt{article\_embedding\_documents}, \texttt{article\_embedding\_chunks};
    \item vector của chunk và category: \texttt{article\_chunk\_embeddings}, \texttt{category\_embedding\_documents}, \texttt{category\_embeddings};
    \item lớp điều phối: \texttt{embedding\_jobs}.
\end{itemize}

Thiết kế này cho phép lưu riêng phần dữ liệu vector mà không làm phình database của Article-Service. Ngoài ra, việc tách model và embedding version giúp hệ thống có thể tái sinh vector khi thay đổi mô hình hoặc chiến lược chunking.

Trong tương lai, nếu kết hợp với \texttt{pgvector}, service này có thể trở thành nền tảng cho tìm kiếm ngữ nghĩa, gợi ý nội dung và hỗ trợ AI sinh tóm tắt hoặc trả lời câu hỏi dựa trên tri thức từ bài báo.

\embeddingservicediagram

\section{Nguyễn Việt Huy -- MSSV: B22DCCN393}

\subsection{Identity-Service}

Identity-Service là nơi quản lý thông tin định danh và xác thực của toàn hệ thống. Đây là service giữ vai trò nền tảng vì hầu hết các service khác đều tham chiếu đến người dùng thông qua external id do service này cung cấp.

Bảng trung tâm là \texttt{users}, nơi lưu email, tên hiển thị, username, trạng thái tài khoản, locale, timezone và các mốc thời gian quan trọng. Từ bảng này, hệ thống mở rộng thành các lớp dữ liệu liên quan đến xác thực:

\begin{itemize}
    \item \texttt{user\_identities}: quản lý các provider đăng nhập như email, Google hoặc Apple;
    \item \texttt{user\_devices}: quản lý thiết bị đã đăng nhập;
    \item \texttt{auth\_sessions}: quản lý refresh token và session theo thiết bị;
    \item \texttt{email\_verification\_tokens}: token xác minh email;
    \item \texttt{password\_reset\_tokens}: token đặt lại mật khẩu.
\end{itemize}

Thiết kế này có ưu điểm là tách biệt rõ profile người dùng và các thông tin xác thực động. Nhờ đó, hệ thống có thể hỗ trợ đa phương thức đăng nhập, theo dõi session trên nhiều thiết bị và kiểm soát an toàn tốt hơn trong các nghiệp vụ như xác minh email hoặc đổi mật khẩu.

\identityservicediagram

\subsection{Notification-Service}

Notification-Service phụ trách lưu thông báo và lịch sử gửi thông báo đến người dùng. Trong bức tranh tổng thể của hệ thống, đây là nơi chuyển đổi các sự kiện nghiệp vụ thành trải nghiệm thông báo cụ thể trong ứng dụng hoặc trên thiết bị.

Các bảng cốt lõi của service gồm:

\begin{itemize}
    \item \texttt{user\_notification\_settings}: lưu cấu hình nhận thông báo của từng người dùng;
    \item \texttt{notifications}: lưu nội dung thông báo hiển thị trong app;
    \item \texttt{delivery\_logs}: lưu log gửi thông báo qua các kênh như push hoặc email;
    \item \texttt{templates}: lưu các mẫu thông báo tái sử dụng.
\end{itemize}

Một điểm hay trong thiết kế này là khả năng lưu snapshot actor và target trong bảng \texttt{notifications}. Nhờ đó, thông báo vẫn có thể hiển thị đúng ngữ cảnh ngay cả khi dữ liệu nguồn ở service khác đã thay đổi. Đây là cách làm phù hợp với microservices vì nó giảm phụ thuộc vào truy vấn đồng bộ sang service khác khi render thông báo.

\notificationservicediagram

\section{Lê Minh Ngọc -- MSSV: 586}

\subsection{Billing-Service}

Billing-Service là service phục vụ các nghiệp vụ tài chính và quyền truy cập nội dung. Đây là domain đòi hỏi tính nhất quán cao nên PostgreSQL là lựa chọn phù hợp nhất.

Schema của Billing-Service được tổ chức thành nhiều lớp:

\begin{itemize}
    \item \texttt{subscription\_plans}: định nghĩa các gói dịch vụ;
    \item \texttt{user\_subscriptions}: lưu đăng ký hiện tại của người dùng;
    \item \texttt{payment\_transactions}: lưu các giao dịch thanh toán;
    \item \texttt{credit\_ledger}: lưu lịch sử biến động credit;
    \item \texttt{credit\_reservations}: lưu trạng thái giữ chỗ credit;
    \item \texttt{content\_entitlements}: lưu quyền truy cập nội dung đã được cấp.
\end{itemize}

Thiết kế này giúp tách rõ ba lớp nghiệp vụ quan trọng: giao dịch tiền tệ, biến động credit và quyền truy cập nội dung. Đây là một ưu điểm lớn vì trong các hệ thống subscription, ba lớp này thường có quan hệ với nhau nhưng không nên gộp chung vào một bảng.

Mặc dù service này hiện chưa được wire đầy đủ vào runtime migration, nhưng schema hiện có cho thấy hướng thiết kế khá chặt chẽ và phù hợp để tích hợp thanh toán ở các giai đoạn sau.

\billingservicediagram

\subsection{Social-Service}

Social-Service được thiết kế theo hướng document store thay vì quan hệ truyền thống. Điều này hợp lý vì nghiệp vụ social thường phát sinh lượng ghi lớn, cấu trúc dữ liệu biến động nhanh và cần tối ưu cho các thao tác feed, comment, reaction hoặc playlist.

Theo tài liệu thiết kế hiện có, Social-Service dự kiến quản lý các collection:

\begin{itemize}
    \item \texttt{user\_follows};
    \item \texttt{show\_follows};
    \item \texttt{saved\_episodes};
    \item \texttt{playlists};
    \item \texttt{playlist\_items};
    \item \texttt{listening\_progress};
    \item \texttt{episode\_reactions};
    \item \texttt{comments};
    \item \texttt{episode\_engagement\_counters}.
\end{itemize}

Mô hình document của Social-Service có một số lợi thế:

\begin{itemize}
    \item linh hoạt khi thay đổi cấu trúc comment hoặc reaction;
    \item dễ lưu snapshot thông tin hiển thị;
    \item phù hợp với workload ghi nhiều và đếm nhiều;
    \item có thể kết hợp với Redis để tối ưu counter hoặc dữ liệu nóng.
\end{itemize}

Vì Social-Service chưa được triển khai đầy đủ trong phần runtime hiện tại, phần này trong báo cáo có thể được xem là định hướng thiết kế hơn là schema đã vận hành thực tế.

\socialservicediagram

\chapter{Kết luận}

Qua quá trình khảo sát và mô tả, có thể thấy cơ sở dữ liệu của hệ thống Pody được xây dựng theo hướng tách domain rõ ràng, phù hợp với kiến trúc microservices và dễ mở rộng trong tương lai. Phần tổng quan toàn hệ thống giúp làm rõ cách các service phối hợp với nhau ở mức dữ liệu, trong khi phần mô tả cá nhân cho phép từng thành viên đi sâu vào các service được phân công.

Việc phân chia nội dung theo cặp service cho từng thành viên cũng khá hợp lý vì:

\begin{itemize}
    \item mỗi người có thể tập trung vào một nhóm domain liên quan;
    \item phần trình bày cá nhân vẫn bám sát tổng thể hệ thống;
    \item báo cáo vừa có chiều rộng ở mức hệ thống, vừa có chiều sâu ở mức service.
\end{itemize}

Trong các bước hoàn thiện tiếp theo, nhóm có thể bổ sung thêm ảnh xuất từ các sơ đồ Mermaid đã xây dựng để thay cho các khung giữ chỗ trong tài liệu này. Khi đó, báo cáo sẽ vừa có phần mô tả bằng lời, vừa có phần trực quan minh họa cho từng service.

\end{document}
